\begin{textsample}{POS Dim 2 – 2014\_c4 – Score 12.00 – 2014\_c4/enade\_2014\_eng\_el\_d9d6525f24.txt}  \label{ex:f2_pos_116}
CAPÍTULO 4 PERCEPÇÃO DA PROVA As análises feitas neste capítulo tratam das percepções dos concluintes da Área de Engenharia Elétrica sobre a prova aplicada no ENADE/2014.
Estas percepções foram mensuradas por meio de nove questões que avaliaram desde o grau de dificuldade da prova até o tempo gasto para concluí -la.
As percepções sobre a prova foram relacionadas com o desempenho dos estudantes e com a Grande Região de funcionamento do curso.
O \textbf{questionário} de percepção da prova encontra -se ao final do Anexo IV, que traz a \textbf{reprodução} da prova.
O desempenho dos estudantes foi classificado em quatro quartos.
Para tanto, esse desempenho foi \textbf{ordenado} de forma ascendente.
O percentil 25, P25, também conhecido como primeiro quartil, é a nota de desempenho que deixa um quarto ( 25 \% ) dos valores observados abaixo e três quartos acima.
A Figura 1 apresenta uma ilustração deste conceito.
O quarto inferior de desempenho é composto pelas notas abaixo do primeiro quartil.
Já o percentil 75, P75, também conhecido como terceiro quartil, é o valor para o qual há três quartos ( 75 \% ) dos dados abaixo e um quarto acima dele.
O quarto superior de desempenho é composto pelas notas iguais ou acima do terceiro quartil.
O percentil 50, P50, também conhecido como mediana, é o valor que divide as notas em dois conjuntos de igual tamanho.
O segundo quarto inclui valores entre o primeiro quartil ( P25 ) e a mediana.
O terceiro quarto contém os valores entre a mediana ( P50 ) e o terceiro quartil ( P75 ).
Vale ressaltar que percentis, quartis e medianas são pontos que não obrigatoriamente pertencem ao conjunto original de dados, ao passo que os quartos são \textbf{subconjuntos} dos dados originais.
A seguir, serão apresentados gráficos com resultados selecionados, relativos às nove questões avaliadas por grupos de estudantes.
Os gráficos apresentam nas barras o percentual de alunos que assinalaram uma das opções ou a soma das porcentagens daqueles que assinalaram duas ( ou três ) delas.
Por exemplo, para as questões 1 e 2, os gráficos apresentam a porcentagem total de participantes que assinalaram as opções ( D ) difícil e ( E ) muito difícil.
Em cada barra foram assinalados também os extremos do intervalo de confiança de 95 \% como linhas verticais unidas por uma linha horizontal na forma da letra H maiúscula.
As Tabelas no Anexo II apresentam os valores absolutos e a distribuição percentual20 das alternativas válidas das nove questões, segundo o mesmo \textbf{recorte} de desempenho dos alunos e Grande Região de funcionamento do curso. 4.1 GRAU DE DIFICULDADE DA PROVA 4.1.1 Componente de Formação Geral Ao avaliarem “ Qual o grau de dificuldade desta prova na parte de Formação Geral? ” ( Questão 1 ), 24,5 \% do grupo de inscritos e presentes optaram pelas alternativas difícil ou muito difícil.
Entretanto, para mais da metade dos estudantes ( 58,8 \% ), o Componente de Formação Geral da prova foi considerado com grau de dificuldade médio ( Gráfico 4.1, Gráfico 4.2 e, no Anexo II, a Tabela II.1 ).
O percentual de estudantes que consideraram a prova como difícil ou muito difícil foi maior na região Centro-Oeste, onde a proporção foi de 29,8 \%, enquanto a de menor incidência foi a Nordeste, com 16,7 \%.
No Gráfico 4.1, é possível observar que esta diferença é estatisticamente significativa.
Nas Grandes Regiões, a proporção de presentes à prova que consideraram o Componente de Formação Geral como sendo de grau de dificuldade médio esteve entre 54,6 \% na região Centro-Oeste e 64,9 \% na região Nordeste. 20Cumpre lembrar uma das convenções para tabelas \textbf{numéricas} ( pág. iii ) sobre a possibilidade da soma das partes não resultar em 100 \% por questões de \textbf{arredondamento}.
O percentual de alunos que consideraram a prova difícil ou muito difícil foi decrescente em função dos quartos de desempenho : 35,2 \% no primeiro quarto e 12,9 \% no quarto, grupo de melhor desempenho na prova.
Nos quartos de desempenho intermediários, a proporção de alunos que consideraram a prova difícil ou muito difícil foi de 28,2 \% no 2º quarto e 21,8 \% no 3º quarto.
Estes valores são decrescentes com o desempenho e as diferenças entre todos os quartos são estatisticamente significativas.
Para todos os quartos de desempenho, a alternativa modal para esta pergunta foi médio, com 54,7 \% e 60,9 \% dos respondentes nos quartos extremos, primeiro e quarto, respectivamente, valores crescentes com o desempenho. 4.1.2 Componente de Conhecimento Específico Ao responderem à Questão 2 – “ Qual o grau de dificuldade desta prova na parte de Componente Específico? ” – 51,6 \% do grupo de estudantes classificaram -na como difícil ou muito difícil.
Além disso, o Componente de Conhecimento Específico da prova foi considerado com grau de dificuldade médio por 44,3 \% dos alunos ( Gráfico 4.3, Gráfico 4.4, e, no Anexo II, a Tabela II.2 ).
A análise das respostas dos estudantes quanto ao grau de dificuldade do Componente de Conhecimento Específico da prova, agregado por Grande Região, mostra que a diferença entre a maior e a menor proporções de alunos que a avaliaram como difícil ou muito difícil é estatisticamente significativa : a maior na região Centro-Oeste ( 57,6 \% ) e a menor proporção na região Nordeste ( 41,0 \% ).
O percentual de alunos que classificaram o grau de dificuldade como médio, no Componente de Conhecimento Específico, variou de 39,1 \% a 54,1 \%, para as regiões Centro-Oeste e Nordeste, respectivamente.
A proporção dos que classificaram a parte específica como difícil ou muito difícil variou de 40,9 \% ( 4º quarto ) a 55,5 \% ( 1º quarto ), valores decrescentes com o desempenho.
Considerando -se a avaliação da dificuldade das questões do Componente de Conhecimento Específico da prova, de acordo com o desempenho dos estudantes, observa -se que há diferença estatisticamente significativa do resultado do quarto de melhor desempenho em relação aos demais quartos.
A alternativa modal para a Questão 2 foi o grau médio, com 40,6 \% do quarto inferior e 52,1 \% do superior optando por esta resposta, valores crescentes com o desempenho. 4.2 EXTENSÃO DA PROVA EM RELAÇÃO AO TEMPO TOTAL \textbf{Indagados} quanto à extensão da prova, em relação ao tempo total oferecido para a sua resolução ( Questão 3 ), os estudantes apontaram, com maior incidência, a alternativa que considerava a extensão adequada, para todas as agregações consideradas ( Gráfico 4.5, Gráfico 4.6, e, no Anexo II, a Tabela II.3 ).
O percentual de alunos que responderam ser a extensão da prova adequada foi de 49,3 \%.
Já 41,5 \% dos inscritos presentes consideraram que o exame foi longo ou muito longo e 9,2 \% o avaliaram como curto ou muito curto.
Entre as Grandes Regiões a proporção daqueles que avaliaram a prova como longa ou muito longa em relação ao tempo total destinado à sua resolução variou de 34,2 \% na região Norte até 48,1 \% na região Sul.
A diferença entre as regiões Norte e Sul é estatisticamente significativa.
Considerando -se o desempenho dos alunos, nota -se ainda que 52,0 \% consideraram a extensão da prova adequada no quarto de desempenho inferior e 44,0 \% no de melhor desempenho ( quarto superior ).
Nos quartos intermediários, esta proporção foi 51,4 \% no segundo quarto e 49,9 \% no terceiro.
Esses valores são decrescentes com o desempenho.
No Gráfico 4.6, pode -se constatar que há, \textbf{grosso} modo, uma tendência crescente da proporção de estudantes que consideraram a prova longa ou muito longa em função dos quartos de desempenho.
Observam -se diferenças estatisticamente significativas entre as proporções de escolha desta resposta no último quarto e nos demais. 4.3 COMPREENSÃO DOS ENUNCIADOS DAS QUESTÕES 4.3.1 Componente de Formação Geral Com relação aos enunciados das questões do Componente de Formação Geral ( Questão 4 ), as opiniões foram positivas, já que 76,4 \% dos alunos avaliados consideraram os enunciados de todas ou da maioria das questões claros e objetivos ( Gráfico 4.7, Gráfico 4.8, e, no Anexo II, a Tabela II.4 ).
Na análise regional, a \textbf{percentagem} de estudantes que avaliaram que todos ou a maioria dos enunciados das questões do Componente de Formação Geral estavam claros e objetivos variou de 71,8 \% na região Sul a 79,4 \% na região Centro-Oeste, sendo esta diferença entre Sul e Centro-Oeste estatisticamente significativa.
A análise das percepções dos estudantes sobre a clareza e objetividade dos enunciados permite afirmar que todos, ou a maioria dos enunciados de questões relativas ao Componente de Formação Geral, foram considerados claros e objetivos para a maior parte dos respondentes ( maior do que 71 \% em todas as regiões e maior do que 68 \% para todos os quartos de desempenho ).
Segundo o desempenho, observa -se que a proporção dos que emitiram esta opinião \textbf{cresce} conforme o desempenho aumenta, com diferenças estatisticamente significativas do primeiro para os demais quartos.
No quarto superior, a clareza e objetividade de todos ou da maioria dos enunciados das questões foi percebida por 80,5 \% dos alunos sem ser estatisticamente significativa a diferença deste para o quarto imediatamente inferior, para o qual a proporção foi de 79,7 \%. 4.3.2 Componente de Conhecimento Específico Com relação aos enunciados das questões do Componente de Conhecimento Específico da prova, para 76,8 \% dos estudantes avaliados da Área de Engenharia Elétrica, a clareza e a objetividade ( Questão 5 ) estavam presentes em todas ou na maioria das questões ( Gráfico 4.9, Gráfico 4.10, e no Anexo II, a Tabela II.5 ).
A maioria dos estudantes de todas as Grandes Regiões brasileiras considerou claros e objetivos todos ou a maioria dos enunciados das questões do Componente de Conhecimento Específico da prova, percentual sempre maior ou igual a 70,4 \%.
A diferença entre a maior e a menor proporção, 79,1 \% ( Nordeste ) e 70,4 \% ( Norte ), é estatisticamente significativa.
A proporção de estudantes que consideraram os enunciados das questões claros e objetivos apresenta uma tendência crescente em relação ao aumento de desempenho : mais elevada no quarto superior ( 83,2 \% ) se comparada ao quarto inferior de desempenho ( 69,5 \% ).
As diferenças entre os quartos são todas estatisticamente significativas, com exceção da diferença entre os dois quartos centrais. 4.4 SUFICIÊNCIA DAS INFORMAÇÕES/INSTRUÇÕES FORNECIDAS Ao avaliarem as informações/instruções fornecidas para a resolução das questões ( Questão 6 ), 84,8 \% dos respondentes da Área de Engenharia Elétrica de todo o Brasil afirmaram que estas eram até excessivas ou suficientes em todas ou na maioria das questões ( Gráfico 4.11, Gráfico 4.12, e, no Anexo II, a Tabela II.6 ).
Quanto à distribuição de respondentes pelas Grandes Regiões observa -se que a proporção de estudantes que consideraram as informações/instruções fornecidas até excessivas ou suficientes em todas ou na maioria das questões foi sempre maior ou igual a 80,1 \%, chegando a 88,4 \% na região Nordeste.
A diferença entre o maior percentual, obtido na região Nordeste ( 88,4 \% ), e o menor, da região Norte ( 80,1 \% ) é estatisticamente significativa.
O percentual de participantes que avaliaram as informações/instruções como até excessivas ou suficientes em todas ou na maioria das questões foi mais elevado no quarto superior ( 91,9 \% ), percentual superior à média nacional ( 84,8 \% ).
Já no quarto inferior, a suficiência das informações/instruções declarada como até excessiva em todas ou na maioria das questões foi percebida por 77,6 \% dos respondentes.
Estas proporções são crescentes com os quartos de desempenho.
Levando -se em conta o desempenho dos participantes, notam -se diferenças estatisticamente significativas entre as opiniões de estudantes de todos os quartos, com exceção dos dois quartos centrais, como mostra o Gráfico 4.12. 4.5 DIFICULDADE ENCONTRADA AO RESPONDER À PROVA Perguntados sobre as dificuldades com as quais se depararam ao responder à prova ( Questão 7 ), 16,5 \% dos estudantes apontaram o desconhecimento do conteúdo.
Para 45,4 \%, a forma diferente de abordagem do conteúdo foi indicada como dificuldade.
Já a falta de motivação para fazer a prova foi a dificuldade apontada por 18,7 \% dos respondentes.
Considerando -se todo o Brasil, 10,9 \% dos respondentes afirmaram que não tiveram qualquer tipo de dificuldade para responder à prova ( Tabela II.7 no Anexo II ).
Os Gráficos 4.13 e 4.14 apresentam os percentuais de estudantes que apontaram o desconhecimento do conteúdo como dificuldade percebida ao responder à prova.
Na análise por Grandes Regiões, o percentual de inscritos e presentes que apontaram o desconhecimento do conteúdo como dificuldade ao responder à prova não superou 21,8 \%.
Os percentuais variaram de 15,7 \% na região Sul a 21,8 \% na Centro-Oeste, sendo esta diferença significativa estatisticamente.
A forma diferente de abordagem do conteúdo foi a escolha modal dos estudantes, com percentuais que variaram de 37,3 \% ( região Nordeste ) a 48,0 \% ( Sudeste ).
O percentual de alunos que citaram a falta de motivação como dificuldade variou de 16,4 \% ( região Sudeste ) a 24,3 \% ( região Nordeste ).
Os que declararam não ter qualquer dificuldade para responder à prova variaram de 7,7 \% na região Centro-Oeste a 11,7 \% na Sudeste.
Com relação aos quartos de desempenho, o desconhecimento do conteúdo foi a opção escolhida por 18,8 \% dos estudantes do quarto superior e 14,7 \% do quarto inferior, valores crescentes com o desempenho.
A alternativa modal para os alunos, quando agregados pelos quartos de desempenho, foi que a dificuldade encontrada foi causada pela forma diferente de abordagem do conteúdo : 47,6 \% no quarto inferior e 36,5 \% do quarto superior assim o responderam. 4.6 CONTEÚDOS DAS QUESTÕES OBJETIVAS DA PROVA Ao analisarem os conteúdos das questões objetivas da prova ( Questão 8 ), um percentual muito pequeno dos estudantes avaliados, apenas 3,7 \%, afirmou que não estudou ainda a maioria desses conteúdos ( Gráficos 4.15, Gráfico 4.16, e a Tabelas II.8 no Anexo II ).
A maioria ( 64,6 \% ) afirmou ter estudado e aprendido muitos ou todos os conteúdos.
Na análise por Grande Região, a proporção de respondentes que escolheram a opção não estudou ainda a maioria desses conteúdos foi também pequena.
Observa -se que nas regiões Centro-Oeste ( 5,8 \% ), Norte ( 3,8 \% ) e Nordeste ( 3,8 \% ), apesar de pequenas, as proporções foram maiores do que a média nacional ( 3,7 \% ).
Observa -se que não há diferença estatisticamente significativa entre a região Norte e as demais regiões.
Em todas as regiões, a maioria dos presentes afirmou ter estudado e aprendido muitos ou todos os conteúdos, com proporções variando entre 53,4 \% na região Centro-Oeste e 69,0 \% na Sul.
Considerando -se separadamente as opiniões de estudantes dos quatro quartos de desempenho, observa -se que, no quarto inferior, 6,4 \% ofereceram como resposta que não estudou ainda a maioria desses conteúdos, sendo 1,9 \% os do quarto superior com a mesma resposta, valores decrescentes com o desempenho.
A diferença entre os alunos que optaram por este motivo de dificuldade nos quartos extremos é estatisticamente significativa.
Tendo em conta o quarto superior, 77,9 \% dos alunos afirmaram ter estudado e aprendido muitos ou todos os conteúdos.
No outro extremo, no primeiro quarto, 53,0 \% optaram pelas mesmas categorias, valores crescentes com o desempenho. 4.7 TEMPO GASTO PARA CONCLUIR A PROVA Ao responderem sobre o tempo de conclusão da prova ( Questão 9 ), mais de dois terços dos estudantes ( 70,0 \% ) afirmou ter gasto entre duas e quatro horas ( Gráfico 4.17, Gráfico 4.18 e, no Anexo II, a Tabela II.9 ).
Considerando -se as cinco Grandes Regiões brasileiras, o percentual dos que utilizaram entre duas e quatro horas para finalizar a prova apenas nas regiões Nordeste ( 65,8 \% ) e Centro-Oeste ( 67,0 \% ) foi inferior ao percentual nacional.
Nas demais Grandes Regiões o percentual de alunos que \textbf{dispensaram} entre duas e quatro horas para concluir a prova, não foi superior a 71,2 \%, como mostra o Gráfico 4.17.
Uma vez sendo analisadas as alternativas escolhidas pelos estudantes que se situam nos diferentes quartos de desempenho, observa -se, \textbf{grosso} modo, uma tendência crescente : uma maior proporção de participantes no quarto superior declarou ter gasto entre duas e quatro horas para concluir a prova quando comparadas com os dos quartos inferiores.
A diferença entre o primeiro ( 64,1 \% ) em relação aos demais quartos de desempenho é estatisticamente significativa.

% matched lemmas: arredondamento, crescer, dispensar, grosso, indagar, numérico, ordenar, percentagem, questionário, recorte, reprodução, subconjunto
\end{textsample}
